% =====================================================================
%   INTERROGATION N° 1  -  MATHÉMATIQUES  -  CLASSE DE 3ème A
%   Leçon : Monômes - Polynômes - Fractions rationnelles
%   Prof : KODJONE Kodjo   -   Établissement : LYCÉE KPETA
%   Format : 1 feuille A4 = 4 exemplaires (2 x version A en haut,
%            2 x version B en bas)  -  impression noir et blanc
%   Page 2 : corrigé (pour le professeur, à ne pas distribuer)
%   Compilation : pdflatex
% =====================================================================
\documentclass[10pt,a4paper]{article}
\usepackage[utf8]{inputenc}
\usepackage[T1]{fontenc}
\usepackage{ae,aecompl}
\usepackage[a4paper,margin=5mm,top=4mm,bottom=4mm]{geometry}
\usepackage{amsmath,amssymb}
\usepackage{array}
\pagestyle{empty}
\setlength{\parindent}{0pt}
\setlength{\fboxsep}{0.6mm}
\setlength{\fboxrule}{0.4pt}
\hyphenpenalty=10000
\exhyphenpenalty=10000
\tolerance=3000

% ---------------------- réglages du tableau QCM ----------------------
\newcommand{\tailleQCM}{\fontsize{7.4}{8.8}\selectfont}
\DeclareMathSizes{7.4}{7.4}{5}{5}
\DeclareMathSizes{6.6}{6.6}{4.6}{4.6}
\DeclareMathSizes{6.4}{6.4}{4.5}{4.5}
\newcolumntype{E}{>{\raggedright\arraybackslash}p{24mm}}
\newcolumntype{P}{>{\centering\arraybackslash}p{18mm}}
\newcolumntype{Q}{>{\centering\arraybackslash}p{9.5mm}}
\newcolumntype{R}[1]{>{\raggedleft\arraybackslash}p{#1}}
\newcommand{\enteteQCM}{%
  \multicolumn{2}{|c|}{} & \multicolumn{3}{c|}{\textbf{Réponses proposées}} & \\ \cline{3-5}
  \textbf{N\textsuperscript{o}} & \textbf{Énoncé} & \textbf{a} & \textbf{b} & \textbf{c} & \textbf{Choix} \\ \hline}

% ---------------------------- version A ------------------------------
\newcommand{\lignesA}{%
1 & Degré du monôme $-5x^4$ & 5 & 4 & 3 & \\ \hline
2 & Réduire $3x^2+5x-7x^2+2x$ & $4x^2+7x$ & $-4x^2+3x$ & $-4x^2+7x$ & \\ \hline
3 & Développer $(2x-3)^2$ & $4x^2-12x+9$ & $4x^2+12x+9$ & $4x^2-9$ & \\ \hline
4 & Développer $(3x-5)(3x+5)$ & $9x^2+25$ & $9x^2-30x+25$ & $9x^2-25$ & \\ \hline
5 & Factoriser $x^2-8x+16$ & $(x+4)^2$ & $(x-4)^2$ & $(x-4)(x+4)$ & \\ \hline
6 & Factoriser $(x+1)(2x-3)$ \mbox{${}+(x+1)(x+5)$} & $(x+1)(3x+2)$ & $(x+1)(x-8)$ & $(x+1)(2x-3)$ & \\ \hline
7 & $P(x)=2x^3-3x+1$ : calculer $P(-1)$ & $-4$ & 6 & 2 & \\ \hline
8 & Condition d'existence de $B$ & $x\neq-2$ & $x\neq2$ & $x\neq2$ et $x\neq-2$ & \\ \hline
9 & Forme simplifiée de $B$ & $\frac{x+2}{x-2}$ & $\frac{x-2}{x+2}$ & $x-2$ & \\ \hline
10 & Valeur de $B$ pour $x=3$ & 5 & $-\frac{1}{5}$ & $\frac{1}{5}$ & \\ \hline
}

% ---------------------------- version B ------------------------------
\newcommand{\lignesB}{%
1 & Degré du monôme $6x^5$ & 6 & 4 & 5 & \\ \hline
2 & Réduire $4x^2+7x-6x^2+2x$ & $2x^2+9x$ & $-2x^2+9x$ & $-2x^2+5x$ & \\ \hline
3 & Développer $(x-5)^2$ & $x^2+10x+25$ & $x^2-25$ & $x^2-10x+25$ & \\ \hline
4 & Développer $(2x-7)(2x+7)$ & $4x^2-49$ & $4x^2+49$ & $4x^2-28x+49$ & \\ \hline
5 & Factoriser $x^2+6x+9$ & $(x-3)^2$ & $(x-3)(x+3)$ & $(x+3)^2$ & \\ \hline
6 & Factoriser $(x-2)(3x+1)$ \mbox{${}-(x-2)(x-5)$} & $(x-2)(4x-4)$ & $(x-2)(2x+6)$ & $(x-2)(3x+1)$ & \\ \hline
7 & $Q(x)=3x^2-2x+5$ : calculer $Q(-1)$ & 10 & 6 & $-4$ & \\ \hline
8 & Condition d'existence de $C$ & $x\neq-3$ & $x\neq3$ et $x\neq-3$ & $x\neq3$ & \\ \hline
9 & Forme simplifiée de $C$ & $\frac{x+3}{x-3}$ & $\frac{x-3}{x+3}$ & $x+3$ & \\ \hline
10 & Valeur de $C$ pour $x=1$ & 2 & $-2$ & $-\frac{1}{2}$ & \\ \hline
}

% ------------------------- un exemplaire complet ----------------------
\newcommand{\unecopie}[3]{%
\fbox{%
\begin{minipage}[t][139.5mm][t]{96.4mm}%
\vspace{0.4mm}%
{\fontsize{6.3}{7.4}\selectfont
\begin{tabular}{@{}p{44mm}@{}R{45mm}@{}}
\textbf{Établissement :} LYCÉE KPETA & Année scolaire : 2026 -- 2027 \\
\textbf{Prof :} KODJONE Kodjo & Classe : 3\textsuperscript{ème} A \\
\end{tabular}}%
\vspace{0.4mm}\hrule\vspace{0.9mm}%
\begin{center}
{\fontsize{8.2}{9.6}\selectfont\textbf{INTERROGATION N\textsuperscript{o} 1 DE MATHÉMATIQUES}}\\[0.6mm]
{\fontsize{6.6}{7.8}\selectfont Leçon : Monômes -- Polynômes -- Fractions rationnelles}\\[0.6mm]
{\fontsize{6.4}{7.4}\selectfont\makebox[96.4mm]{Durée : 30 min \hfill Note : \dots\dots\dots\,/\,10 \hfill \textbf{VERSION #1}}}
\end{center}
\vspace{0.4mm}%
{\fontsize{6.6}{7.8}\selectfont Nom : \dotfill\ Prénom : \dotfill\ Classe : 3\textsuperscript{ème} A}%
\vspace{0.9mm}%
\begin{center}
{\fontsize{6.6}{7.8}\selectfont \textbf{Consigne :} Indique la bonne réponse par \textbf{a}, \textbf{b} ou \textbf{c}.}\\[0.5mm]
{\fontsize{6.4}{7.6}\selectfont #2}
\end{center}
\vspace{0.4mm}%
\begingroup
\setlength{\tabcolsep}{0.7pt}%
\renewcommand{\arraystretch}{1.4}%
\arrayrulewidth=0.4pt%
\tailleQCM
\begin{tabular}{|c|E|P|P|P|Q|}
\hline
\enteteQCM
#3%
\end{tabular}
\endgroup
\end{minipage}}%
}

\newcommand{\noteA}{Questions 8 à 10 : on pose $B=\frac{x^{2}-4}{x^{2}+4x+4}$.}
\newcommand{\noteB}{Questions 8 à 10 : on pose $C=\frac{x^{2}-9}{x^{2}-6x+9}$.}

% =====================================================================
\begin{document}

% ------------- FEUILLE ÉLÈVE : 4 EXEMPLAIRES SUR UNE FEUILLE A4 ------
\hfill\unecopie{A}{\noteA}{\lignesA}\hspace{1.6mm}\vrule width 0.3pt\hspace{1.6mm}\unecopie{A}{\noteA}{\lignesA}\hfill

\vspace{1mm}
\noindent\hrule height 0.3pt
\vspace{1mm}

\hfill\unecopie{B}{\noteB}{\lignesB}\hspace{1.6mm}\vrule width 0.3pt\hspace{1.6mm}\unecopie{B}{\noteB}{\lignesB}\hfill

\newpage

% ------------- PAGE DE CORRECTION (PROFESSEUR) -----------------------
\begin{center}
{\Large\textbf{CORRIGÉ -- INTERROGATION N\textsuperscript{o} 1 DE MATHÉMATIQUES}}\\[2mm]
{\normalsize Lycée Kpeta --- Classe de 3\textsuperscript{ème} A --- Prof : KODJONE Kodjo --- Année scolaire 2026 -- 2027}\\[1.5mm]
{\normalsize Leçon : Monômes -- Polynômes -- Fractions rationnelles --- Durée : 30 min --- Note sur 10}\\[1.5mm]
{\small 1 point par réponse juste ; aucun point négatif.}
\end{center}

\vspace{3mm}
\noindent{\small
\begin{tabular}{|c|c|c|p{112mm}|}
\hline
\textbf{N\textsuperscript{o}} & \textbf{Version A} & \textbf{Version B} & \textbf{Justification} \\ \hline
1 & \textbf{b} & \textbf{c} & A : monôme $-5x^{4}$, le degré est l'exposant de $x$ : 4. \quad B : $6x^{5}$, degré 5. \\ \hline
2 & \textbf{c} & \textbf{b} & A : $3x^{2}-7x^{2}=-4x^{2}$ et $5x+2x=7x$, donc $-4x^{2}+7x$. \quad B : $4x^{2}-6x^{2}=-2x^{2}$ et $7x+2x=9x$, donc $-2x^{2}+9x$. \\ \hline
3 & \textbf{a} & \textbf{c} & A : $(2x-3)^{2}=(2x)^{2}-2\times2x\times3+3^{2}=4x^{2}-12x+9$. \quad B : $(x-5)^{2}=x^{2}-10x+25$. \\ \hline
4 & \textbf{c} & \textbf{a} & A : $(3x-5)(3x+5)=(3x)^{2}-5^{2}=9x^{2}-25$. \quad B : $(2x-7)(2x+7)=(2x)^{2}-7^{2}=4x^{2}-49$. \\ \hline
5 & \textbf{b} & \textbf{c} & A : $x^{2}-8x+16=x^{2}-2\times x\times4+4^{2}=(x-4)^{2}$. \quad B : $x^{2}+6x+9=x^{2}+2\times x\times3+3^{2}=(x+3)^{2}$. \\ \hline
6 & \textbf{a} & \textbf{b} & A : $(x+1)\left[(2x-3)+(x+5)\right]=(x+1)(3x+2)$. \quad B : $(x-2)\left[(3x+1)-(x-5)\right]=(x-2)(2x+6)$. \\ \hline
7 & \textbf{c} & \textbf{a} & A : $P(-1)=2(-1)^{3}-3(-1)+1=-2+3+1=2$. \quad B : $Q(-1)=3(-1)^{2}-2(-1)+5=3+2+5=10$. \\ \hline
8 & \textbf{a} & \textbf{c} & A : $x^{2}+4x+4=(x+2)^{2}$ s'annule pour $x=-2$, donc $x\neq-2$. \quad B : $x^{2}-6x+9=(x-3)^{2}$ s'annule pour $x=3$, donc $x\neq3$. \\ \hline
9 & \textbf{b} & \textbf{a} & A : $B=\frac{(x-2)(x+2)}{(x+2)^{2}}=\frac{x-2}{x+2}$ pour $x\neq-2$. \quad B : $C=\frac{(x-3)(x+3)}{(x-3)^{2}}=\frac{x+3}{x-3}$ pour $x\neq3$. \\ \hline
10 & \textbf{c} & \textbf{b} & A : $B(3)=\frac{3-2}{3+2}=\frac{1}{5}$. \quad B : $C(1)=\frac{1+3}{1-3}=\frac{4}{-2}=-2$. \\ \hline
\end{tabular}

\vspace{4mm}
\noindent\textbf{Rappel de correction :} aux questions 8 à 10, la condition d'existence doit être donnée avant toute simplification ; accepter la réponse exacte quelle que soit la justification écrite par l'élève.
}

\end{document}
